\section{与巨头竞争：四个不与时间窗口}

上一章讲产品定义，本章讲竞争。祝铭明引用马云总结的创业公司与巨头竞争的四个机会：看不见、看不懂、看不上、来不及。智能眼镜这件事，大公司一定看得见，也大概率看得懂；创业公司的窗口主要来自“看不上”到“来不及”之间的两三年。Rokid 的目标不是与巨头平等竞争，而是做到对等竞争，至少上“大人那一桌吃饭”。

\begin{figure}[H]
\centering
\includegraphics[width=0.96\textwidth]{startup-vs-giants-four-no.png}
\caption{创业公司对巨头的四个不：看不见、看不懂、看不上、来不及。自制概念图，依据 01:19:29--01:23:38 对谈内容整理。}
\end{figure}

\begin{knowledgebox}{读图：Rokid 的窗口不是巨头不知道，而是巨头尚未全力}
图中“来不及”是创业公司真正想争取的状态。智能眼镜足够显眼，大厂不会长期看不见；Rokid 必须在大厂全面进入前，通过产品、用户、开发者和供应链形成足够领先。
\end{knowledgebox}

\begin{knowledgebox}{术语消化：四个不}
\noindent\begin{tabularx}{\textwidth}{p{0.16\textwidth}p{0.36\textwidth}X}
\toprule
阶段 & 巨头状态 & 创业公司的窗口 \\
\midrule
看不见 & 机会太边缘，组织没有注意到 & 小团队可以先验证新需求。 \\
看不懂 & 看到了现象，但无法判断价值 & 创业公司用产品和数据证明意义。 \\
看不上 & 看懂了，但规模暂时不够大 & 这是智能眼镜当前最可能的窗口。 \\
来不及 & 等巨头下场时，先发者已有体验和生态积累 & 创业公司要争取进入这个状态。 \\
\bottomrule
\end{tabularx}
\end{knowledgebox}

\subsection{大厂为什么会晚一步}

祝铭明认为，大公司通常敢为人后，选择成熟供应链、成熟方案和更低风险切入。例如不带显示的眼镜更容易用已有供应链做出来；但带显示、轻便、all-in-one 的智能眼镜需要探索新工艺，创业公司反而可能先跑到最前面。问题是，先跑意味着要替行业承担工艺、产能和用户教育风险。

\begin{warningbox}{先发不是单纯优势}
先发者要承担未验证工艺、供应链爬坡、产品缺陷被放大、用户预期管理和资金压力。大厂晚一步进入，反而可以复制成熟路线；创业公司必须用时间差换取口碑、开发者和产品体验。
\end{warningbox}

\subsection{对等竞争：不是平等竞争}

本节解释祝铭明说的“上大人那一桌吃饭”。创业公司无法和腾讯、阿里、字节、小米、苹果、Meta 做平等竞争，因为资金、流量、供应链和渠道都不在同一量级。能争取的是对等竞争：在一个具体产品定义上，用户愿意把 Rokid 和巨头放在同一个选择集合里比较；渠道愿意认真卖；开发者愿意适配；模型和内容伙伴愿意合作。

访谈里提到几个衡量信号：日均使用超过两小时、开发者超过一万、企业开发者超过三千、周活跃超过七成、复购和推荐占比较高。这些数字不应被当作外部审计后的结论，而应被当作讲者用来说明“平台属性正在形成”的证据类型。真正重要的是：创业公司不能只靠发布会热度上桌，必须拿出用户时长、复购、开发者和供应链交付。

\noindent\begin{tabularx}{\textwidth}{p{0.20\textwidth}p{0.32\textwidth}X}
\toprule
对等竞争指标 & 为什么重要 & 不足之处 \\
\midrule
用户时长 & 说明设备不是一次性新鲜感 & 还要看任务结构和留存。 \\
复购/推荐 & 说明早期用户有口碑扩散 & 样本可能偏极客，需要扩大到普通用户。 \\
开发者数量 & 说明有生态供给苗头 & 开发者活跃度和收入仍要验证。 \\
渠道复进货 & 说明产品能在真实销售里滚动 & 首批铺货不等于渠道扎实。 \\
供应链产能 & 说明热度能转成交付 & 高产能也会放大质量问题。 \\
\bottomrule
\end{tabularx}

\subsection{生态与盟军}

Rokid 不认为自己可以单独建完整生态。祝铭明提到，腾讯有内容、社交和场景，豆包、通义千问、DeepSeek 等模型能力也可以接入。硬件创业公司需要把平台、技术和产品体验做好，但生态建设需要更多盟军。这里的现实判断是：智能眼镜如果要成为入口，不可能只靠硬件参数取胜，还要有内容、应用、开发者、模型和渠道共同参与。

\begin{importantbox}{生态不是单家公司自建出来的}
Rokid 可以先承担探索和领先，把技术、平台和产品体验打扎实；但内容、社交、模型、开发者和场景伙伴需要共同进入。对随身 AI 入口来说，硬件公司如果既想做硬件、OS、模型、内容、应用和渠道，资源消耗会迅速超过创业公司承受范围。
\end{importantbox}

\subsection{本章小结}

智能眼镜是巨头一定会进入的战场。创业公司的窗口在于更早承担风险、更快形成产品体验和生态位置，并在巨头全力进入前建立对等竞争资格。

